\begin{pastexam}[label={px:nchu-114-dsa-19}]{中興資工所甲組 114 學年度 單選題第 19 題}
（單選；每題 3 分，未答題得 0 分，答錯不倒扣）

In an array-based circular queue, how do you check if the queue is full?

(A)~By comparing the front and rear indices directly\\
(B)~By checking if (rear + 1) \% size == front\\
(C)~By ensuring rear == front at all times\\
(D)~By ensuring (front + 1) \% size == rear
\end{pastexam}

\begin{pastsoln}
答案是 (B)。

常見的約定（資結書第~6.2.2~節）：陣列 $Q[0..\mathit{size} - 1]$，$\mathit{front}$ 是前端元素的位置，
$\mathit{rear}$ 是下一個要放的位置。 只看兩個索引時，$\mathit{front} = \mathit{rear}$ 已經用來表示空，
滿的時候不能也是這樣，所以留一格不用：尾端再往後一格就撞到前端時算滿，條件是
$(\mathit{rear} + 1) \bmod \mathit{size} = \mathit{front}$，最多放 $\mathit{size} - 1$ 個。
\begin{itemize}
\item (A) 直接比較 $\mathit{front}$ 和 $\mathit{rear}$（相等）是判斷空的條件；不留空格的話空和滿都是相等，分不出來。 錯。
\item (B) 對，見上。
\item (C) 「rear == front at all times」不是一個判斷條件，而且 $\mathit{front} = \mathit{rear}$ 表示空。 錯。
\item (D) $(\mathit{front} + 1) \bmod \mathit{size} = \mathit{rear}$ 表示佇列裡剛好只有一個元素。 錯。
\end{itemize}
另外兩種做法是多記元素個數 $\mathit{count}$（$\mathit{count} = \mathit{size}$ 為滿）或多記一個旗標，就可以用滿
全部的格子；但只有 (B) 是正確的判斷式。
\end{pastsoln}
